\paragraph{C1: Development cohort size.}
We fit GraphSAGE and graph-code boosting on 25\%, 50\%, and 100\% of development patients and held the test cohort fixed. Smaller cohorts retain all eligible landmarks from each sampled patient. Each seed draws its own patient subset, so variation across seeds includes the training sample. Table~\ref{tab:scale-ap} summarizes the learning curves. Both model families are trained on the same seeded patient subsets and scored on the same test landmarks.

\begin{table}[t]
\centering
\caption{Test AUROC and AUPRC by development cohort fraction (mean\(\pm\)SD over five seeds).}
\label{tab:scale-ap}
\small
\begin{tabular}{lcccccc}
\hline
 & \multicolumn{2}{c}{CM} & \multicolumn{2}{c}{Heme} & \multicolumn{2}{c}{MH} \\
Model and cohort & AUROC & AUPRC & AUROC & AUPRC & AUROC & AUPRC \\
\hline
GraphSAGE, 25\% & 0.542\(\pm\)0.021 & 0.137\(\pm\)0.020 & 0.547\(\pm\)0.035 & 0.081\(\pm\)0.025 & 0.540\(\pm\)0.012 & 0.073\(\pm\)0.003 \\
GraphSAGE, 50\% & 0.550\(\pm\)0.030 & 0.131\(\pm\)0.015 & 0.612\(\pm\)0.029 & 0.088\(\pm\)0.017 & 0.517\(\pm\)0.036 & 0.069\(\pm\)0.009 \\
GraphSAGE, 100\% & 0.614\(\pm\)0.018 & 0.160\(\pm\)0.011 & 0.654\(\pm\)0.018 & 0.117\(\pm\)0.004 & 0.492\(\pm\)0.019 & 0.063\(\pm\)0.008 \\
Boosting, 25\% & 0.573\(\pm\)0.039 & 0.137\(\pm\)0.026 & 0.563\(\pm\)0.011 & 0.071\(\pm\)0.004 & 0.572\(\pm\)0.025 & 0.090\(\pm\)0.021 \\
Boosting, 50\% & 0.558\(\pm\)0.044 & 0.152\(\pm\)0.033 & 0.583\(\pm\)0.035 & 0.079\(\pm\)0.014 & 0.558\(\pm\)0.048 & 0.076\(\pm\)0.011 \\
Boosting, 100\% & 0.535\(\pm\)0.000 & 0.138\(\pm\)0.000 & 0.605\(\pm\)0.000 & 0.093\(\pm\)0.000 & 0.611\(\pm\)0.000 & 0.089\(\pm\)0.000 \\

\hline
\end{tabular}
\end{table}
The full-minus-quarter AUROC differences for GraphSAGE were +0.072, +0.107, -0.048 (CM, Heme, MH). The full-minus-quarter AUROC differences for boosting were -0.038, +0.042, +0.040 (CM, Heme, MH). Both families were trained on the same seeded patient subsets.


\paragraph{D1: Edge topology.}
We replaced Layer~U with five degree-preserving rewires and a degree-unconstrained edge shuffle. The GraphSAGE architecture, static node inputs, and evaluation protocol were fixed. Since the replacement graphs retain diagnosis identities and features computed on Layer~U, this estimates the contribution of message-passing edges conditional on those features.

\begin{table}[t]
\centering
\caption{Test AUROC and AUPRC for observed and randomized edges (GraphSAGE). Rewires report the range of the five graphs' seed means.}
\label{tab:null-ap}
\small
\begin{tabular}{lcccccc}
\hline
 & \multicolumn{2}{c}{CM} & \multicolumn{2}{c}{Heme} & \multicolumn{2}{c}{MH} \\
Edges & AUROC & AUPRC & AUROC & AUPRC & AUROC & AUPRC \\
\hline
Layer U & 0.160\(\pm\)0.011 & 0.117\(\pm\)0.004 & 0.063\(\pm\)0.008 \\
Degree-preserving rewires & 0.140--0.171 & 0.084--0.122 & 0.061--0.076 \\
Edge shuffle & 0.155\(\pm\)0.015 & 0.113\(\pm\)0.005 & 0.074\(\pm\)0.004 \\

\hline
\end{tabular}
\end{table}
For CM, the observed-edge mean (0.160) was within the rewire range (0.140--0.171), and the shuffled-edge mean was 0.155. For Heme, the observed-edge mean (0.117) was within the rewire range (0.084--0.122), and the shuffled-edge mean was 0.113. For MH, the observed-edge mean (0.063) was within the rewire range (0.061--0.076), and the shuffled-edge mean was 0.074. The observed graph remained within the rewire range on every target. The learned scores therefore provide limited evidence of an advantage from the particular Layer~U edges.

